\section{Résultats}
\label{sec:results}

\subsection{Parité FP32 PC$\leftrightarrow$carte (mesuré sur carte)}

En FP32, la parité prédiction-à-prédiction est \emph{exacte} en mode gelé (\num{1.000}) sur les deux jeux,
et quasi exacte en mode en ligne (\num{0.9754} Pronostia, \num{0.9887} Monitoring) --- les écarts en ligne
proviennent du \texttt{float32} de la carte contre le \texttt{float64} du PC sur les frontières de
décision. Le Tableau~\ref{tab:parity} récapitule ; la Figure~\ref{fig:parity} l'illustre.

\begin{table}[t]
\centering
\caption{Parité FP32 PC$\leftrightarrow$carte et performance (\emph{mesuré sur carte}, \texttt{exp\_S36}).}
\label{tab:parity}
\begin{tabular}{lccccc}
\toprule
Jeu & Cond. & F1 (PC/carte) & Parité gelé & Parité en ligne \\
\midrule
Pronostia  & 5feat & \num{0.9164} & \num{1.000} & \num{0.9754} \\
Pronostia  & all   & \num{0.9180} & \num{1.000} & \num{0.9626} \\
Monitoring & 5feat & \num{0.9194} & \num{1.000} & \num{0.9887} \\
\bottomrule
\end{tabular}
\end{table}

\begin{figure}[t]
\centering
\includegraphics[width=0.9\linewidth]{fig1_parity_fp32_pc_board.png}
\caption{Parité FP32 PC$\leftrightarrow$carte : gelé (\num{1.000}) vs en ligne. \emph{Mesuré sur carte}
(\texttt{exp\_S36\_parity\_*.json}).}
\label{fig:parity}
\end{figure}

\paragraph{Latence et RAM (Gap~2/Gap~3).} Toutes les latences carte restent trois ordres de grandeur sous
le budget : inférence \SI{50}{\micro\second} (Pronostia) / \SI{48}{\micro\second} (Monitoring) ;
inférence + mise à jour continue \SI{251}{\micro\second} / \SI{239}{\micro\second}, contre
$\approx$\,\SIrange{33}{35}{\micro\second} pour l'inférence PC de référence. L'empreinte \texttt{.bss} varie
de \SIrange{100}{145}{\kilo\byte}, sous les \SI{256}{\kilo\byte} disponibles. La
Figure~\ref{fig:latency} situe ces valeurs face à la ligne des \SI{100}{\milli\second}.

\begin{figure}[t]
\centering
\includegraphics[width=0.9\linewidth]{fig2_latency_gap2.png}
\caption{Latences carte (inférence vs inférence + mise à jour) contre le budget Gap~2 (\SI{100}{\milli
\second}). \emph{Mesuré sur carte} (\texttt{exp\_S36\_summary.json}).}
\label{fig:latency}
\end{figure}

\subsection{Effondrement de la PTQ naïve (mesuré sur carte)}

La quantification INT8 \emph{legacy} de la tête EWC binaire s'effondre sur la carte : le F1 chute de
$\approx$\,\num{0.92} (FP32) à \num{0.138} (Pronostia) et \num{0.1337} (Monitoring), pour un accord
INT8$\leftrightarrow$FP32 de seulement $\approx$\,\numrange{0.60}{0.74}. La RAM des poids est bien divisée
par~4, mais la métrique est détruite (Tableau~\ref{tab:int8board}).

\begin{table}[t]
\centering
\caption{INT8 \emph{legacy} vs FP32 sur carte (\emph{mesuré}, \texttt{exp\_S36}). RAM $\div 4$ mais
F1 effondré.}
\label{tab:int8board}
\begin{tabular}{lcccc}
\toprule
Jeu & F1 FP32 & F1 INT8 legacy & Accord INT8/FP32 & RAM $\div$ \\
\midrule
Pronostia  & \num{0.9164} & \num{0.138}  & \num{0.7364} & $\times 4$ \\
Monitoring & \num{0.9194} & \num{0.1337} & \num{0.5955} & $\times 4$ \\
\bottomrule
\end{tabular}
\end{table}

\subsection{Diagnostic et récupération (émulé PC bit-exact)}

Une ablation \emph{émulée bit-à-bit sur PC} (référence exacte du kernel) isole la cause et démontre la
récupération (Tableau~\ref{tab:ablation}, Figure~\ref{fig:ablation}). Le facteur dominant est
l'\emph{échelle par-tenseur calibrée} sur accumulateur \texttt{int32} : \texttt{per\_tensor\_calib}
restaure $+\,\num{0.88}$ de F1 (Pronostia \num{0.9462}, Monitoring \num{0.9201}). Le format \texttt{q15}
égale exactement le FP32.

\begin{table}[t]
\centering
\caption{Échelle d'ablation INT8 (F1, \emph{émulé PC bit-exact}, \texttt{exp\_S39\_ablation}).}
\label{tab:ablation}
\begin{tabular}{lcc}
\toprule
Schéma & Pronostia F1 & Monitoring F1 \\
\midrule
FP32 (réf.)        & \num{0.9616} & \num{0.9194} \\
\texttt{legacy\_c} & \num{0.0663} & \num{0.1178} \\
\texttt{per\_tensor\_calib} & \textbf{\num{0.9462}} & \textbf{\num{0.9201}} \\
\texttt{per\_channel\_int8} & \num{0.9426} & \num{0.9187} \\
\texttt{q15}       & \num{0.9616} & \num{0.9194} \\
\bottomrule
\end{tabular}
\end{table}

\begin{figure}[t]
\centering
\includegraphics[width=0.9\linewidth]{fig3_ablation_ladder.png}
\caption{Échelle d'ablation \texttt{legacy\_c}$\to$\texttt{per\_tensor\_calib}$\to$\texttt{per\_channel}
$\to$\texttt{q15}. \emph{Émulé PC bit-exact} (\texttt{exp\_S39\_ablation}).}
\label{fig:ablation}
\end{figure}

\begin{figure}[t]
\centering
\includegraphics[width=0.9\linewidth]{fig4_int8_recovery_board.png}
\caption{INT8 vs FP32 : effondrement \emph{legacy} puis récupération. Effondrement et récupération sont
tous deux \emph{mesurés sur carte} (Sprints~36 et~40) ; la colonne per-canal \emph{émulée PC} (Sprint~39)
est rappelée en hachuré. Le format \texttt{q15} sur carte reste « à mesurer ».}
\label{fig:recovery}
\end{figure}

\subsection{Récupération mesurée sur carte (kernel v2 calibré)}
\label{sec:recovery-board}

La récupération n'est plus une prédiction d'émulateur : le kernel v2 à échelle calibrée a été flashé sur
la NUCLEO-F439ZI et streamé sur les quatre cellules per-canal (deux jeux $\times$ deux protocoles). Le F1
revient au niveau FP32 --- \num{0.9173} contre \num{0.9194} en FP32 sur Monitoring, \num{0.8995} contre
\num{0.9164} sur Pronostia --- alors que la RAM des poids reste divisée par~4. La parité vis-à-vis de
l'émulateur bit-exact est \emph{exacte} en mode gelé (\num{1.000}), ce qui valide le portage sans perte, et
aucune erreur CRC n'a été observée (Tableau~\ref{tab:recovery}).

\begin{table}[t]
\centering
\caption{Récupération INT8 \emph{mesurée sur carte} (kernel v2 per-canal, \texttt{exp\_S40\_board\_v2}).
Parité = accord avec l'émulateur bit-exact ; accord = INT8 contre FP32 sur les mêmes échantillons.}
\label{tab:recovery}
\begin{tabular}{llccccc}
\toprule
Jeu & Protocole & F1 & Parité & Accord & Latence (\micro s) & \texttt{.bss} (B) \\
\midrule
Monitoring & gelé     & \num{0.9173} & \num{1.000}  & \num{0.9996} & 65  & 101236 \\
Monitoring & en ligne & \num{0.9016} & \num{0.9885} & \num{0.8337} & 577 & 101236 \\
Pronostia  & gelé     & \num{0.8995} & \num{1.000}  & \num{0.9951} & 68  & 106152 \\
Pronostia  & en ligne & \num{0.9212} & \num{0.9736} & \num{0.8100} & 606 & 106152 \\
\bottomrule
\end{tabular}
\end{table}

Deux nuances méritent d'être dites. D'une part l'accord INT8$\leftrightarrow$FP32 chute en mode en ligne
(\num{0.8337} et \num{0.8100} contre $>\,\num{0.99}$ en gelé) : la mise à jour continue fait diverger les
deux trajectoires de poids, sans que la métrique finale en souffre. D'autre part le format \texttt{q15} et
la reprise A/B du schéma \texttt{int8\_legacy} sous le kernel v2 n'ont pas été streamés : ces cellules
restent « à mesurer », conformément à la règle « aucun chiffre inventé ».

\subsection{Compromis RAM $\times$ F1 $\times$ latence}

La Figure~\ref{fig:pareto} synthétise le compromis : le kernel calibré ramène la métrique au niveau FP32
tout en conservant le gain RAM (INT8 $\div 4$, Q15 $\div 2$), sans coût de latence prohibitif. L'étoile
marque le point réellement flashé (kernel v2 per-canal) ; les autres points sont émulés. Le gain de latence
espéré, lui, n'a pas lieu --- nous y revenons en Section~\ref{sec:discussion}.

\begin{figure}[t]
\centering
\includegraphics[width=0.9\linewidth]{fig5_pareto_ram_f1_latency.png}
\caption{Front de compromis RAM $\times$ F1 $\times$ latence. Sources mixtes \emph{mesuré}/\emph{émulé}
annotées par point.}
\label{fig:pareto}
\end{figure}
